% Shared parent preamble for the PhD thesis.
% Loaded by main.tex, and by every chapters/*.tex subfile via
% \documentclass[../main.tex]{subfiles}.

% --- Core math ---
\usepackage{amsmath,amssymb,amsfonts,mathtools,mathrsfs,dsfont,cancel,bm}
\usepackage{upgreek}
\usepackage{nicefrac}
\usepackage{commath}
\usepackage{braket}
\usepackage[nointegrals]{wasysym}
\usepackage{pifont}
\usepackage{ar}           % supplies \AR, the aspect ratio N/f in internal_tide_eddy
% Italic capital Xi, used in internal_tide_eddy for the placeholder field whose
% flux fraction is computed. No package loaded here defines it, so it is defined
% below. \providecommand steps aside if some future package supplies its own.
\providecommand{\varXi}{\mathit{\Xi}}

% --- Figures, tables ---
\usepackage{graphicx}

\usepackage[export]{adjustbox}
\usepackage{subcaption}
\usepackage{wrapfig}
\usepackage{float}
\usepackage{booktabs}
\usepackage{longtable}  % the list of symbols runs over several pages
\usepackage{dcolumn}
\usepackage{tikz,pgf,import}
\DeclareCaptionJustification{justified}{\justifying}
\captionsetup{justification=justified,singlelinecheck=false,labelfont=large}
\captionsetup[subfigure]{position=top,belowskip=0pt,justification=centering}

% --- Colors, highlighting, track changes ---
\usepackage[dvipsnames]{xcolor}
\usepackage{soul}
\usepackage[normalem]{ulem}
% \usepackage[inline]{trackchanges}

% --- Layout ---
% Wide text block, as in the reference thesis: the default book-class
% margins leave the text only ~121mm wide on A4, this widens it to ~160mm.
% headheight holds two lines, so a long chapter title in the running header
% cannot run into the rule above the text.
\usepackage[a4paper,inner=25mm,outer=25mm,top=25mm,bottom=30mm,headheight=26pt,headsep=12pt]{geometry}
\usepackage{ragged2e}
% No hyphenation anywhere. 10000 is TeX's infinity, so no word is ever split.
\hyphenpenalty=10000
\exhyphenpenalty=10000
\sloppy
\usepackage{parskip}
\usepackage{setspace}
\usepackage{microtype}
\usepackage{enumitem}

\setstretch{1.25}
\setlength{\parskip}{1.25em}
% --- Hyperlinks ---
% Loaded here, before orcidlink, because orcidlink silently \RequirePackages
% hyperref with no options if it isn't already loaded -- loading it later
% with options triggers "Option clash for package hyperref".
\usepackage[
    colorlinks=true,
    linkcolor=blue!60!black,
    citecolor=blue!60!black,
    urlcolor=blue!60!black
]{hyperref}
% \usepackage[nameinlink]{cleveref}

\usepackage{orcidlink}
\usepackage{subfiles}
\usepackage{xspace}  % used by the internal_tide_eddy chapter's \lr/\ir/\hr-family macros

% --- Bibliography ---
% One combined reference list at the end of the thesis. Resources are
% searched in the order listed, so the shared Zotero file is tried first
% and each paper's own .bib is the fallback for keys that predate it.
% APS-style citations: bracketed numbers, numbered in order of first use,
% consecutive runs compressed to [1-3].
\usepackage[backend=biber,style=numeric-comp,sorting=none,natbib=true]{biblatex}
\addbibresource{/Users/rajarshi/Dropbox/LaTeX/global_references_zotero.bib}
% Copies of each paper's own .bib, kept here because the paper folders
% themselves are not uploaded.
\addbibresource{chapters/bib_active_caustics.bib}
\addbibresource{chapters/bib_internal_tide_eddy.bib}
\addbibresource{chapters/bib_thermalization.bib}
\addbibresource{chapters/bib_inertial_caustics.bib}
\addbibresource{chapters/bib_polydisperse.bib}

% The entry title carries the link: to the DOI when the entry has one,
% otherwise to the URL. The DOI, URL and arXiv eprint are never printed
% as standalone fields.
\ExecuteBibliographyOptions{doi=false,url=false,eprint=false,isbn=false}
\newbibmacro*{linkedtitle}[1]{%
  \iffieldundef{doi}
    {\iffieldundef{url}
       {#1}
       {\href{\thefield{url}}{#1}}}
    {\href{https://doi.org/\thefield{doi}}{#1}}}
\DeclareFieldFormat{title}{\usebibmacro{linkedtitle}{\mkbibemph{#1}}}
\DeclareFieldFormat[article,inbook,incollection,inproceedings,patent,thesis,unpublished]%
  {title}{\usebibmacro{linkedtitle}{\mkbibquote{#1\isdot}}}

% --- Thesis accent color + chapter/section style ---
% Matches the grey "Chapter N" numeral with a bold title set below it,
% used as the visual reference for this thesis.
\definecolor{thesisaccent}{HTML}{8C3A2E}
\usepackage{titlesec}
% Chapter heads are set in TeX Gyre Adventor (family qag), a geometric sans
% whose round letters are much wider than the default one, at sizes a few
% steps above \Huge. Everything else keeps the ordinary sans.
\titleformat{\chapter}[display]
  {\normalfont\fontfamily{qag}\fontsize{36}{48}\selectfont\color{gray}}
  {\chaptertitlename\ \thechapter}
  {10pt}
  {\raggedleft\normalfont\sffamily\bfseries\Huge\color{black}}
\titleformat{\section}
  {\normalfont\sffamily\bfseries\fontsize{18}{20}\selectfont}
  {\thesection}{1em}{}
\titleformat{\subsection}
  {\normalfont\sffamily\bfseries\fontsize{14}{16}\selectfont}
  {\thesubsection}{0.75em}{}
\titleformat{\subsubsection}
  {\normalfont\sffamily\bfseries\itshape}
  {\thesubsubsection}{0.5em}{}

% --- Chapter summary box ---
% Opens each chapter in place of the paper's abstract. breakable lets a long
% summary run across a page break.
\usepackage[most]{tcolorbox}
\definecolor{clbox}{RGB}{100, 87, 166}
\tcbset{colframe = clbox}
\newtcolorbox{summary}{%
  breakable, enhanced, sharp corners,
  boxrule=0pt, colframe=clbox!6,
  borderline west={3pt}{0pt}{clbox},
  colback=clbox!6, colbacktitle=clbox!6, coltitle=clbox,
  title=Summary, fonttitle=\rmfamily\bfseries,
  fontupper=\small\itshape
}

% --- Running headers ---
\usepackage{fancyhdr}
\pagestyle{fancy}
\fancyhf{}
\fancyhead[LE]{\thepage\quad\nouppercase{\leftmark}}
\fancyhead[RO]{\nouppercase{\rightmark}\quad\thepage}
\renewcommand{\headrulewidth}{0.4pt}

\setlength{\parskip}{0.8em}
\setlength{\parindent}{0pt}

% ====================== Shared macros ======================
% Merged from Inertial_caustics/r_preamble.tex, polydisperse/r_preamble.tex,
% and the inline preambles of the Active Caustics, JPO and Thermalization
% manuscripts. Where two papers defined the same name differently, the
% more common / more recent version was kept.

\newcommand{\Dp}[3]{\frac{\partial^{#3} #1}{\partial #2^{#3}}}
\newcommand{\Dpc}[3]{\frac{\cancel{\partial}^{#3}#1}{\cancel{\partial} #2^{#3}}}
\newcommand{\dk}[2]{\frac{d^{#2} #1}{(2\pi)^{#2}}}
\renewcommand{\bra}[1]{\left<#1\right|}
\renewcommand{\ket}[1]{\left|#1\right>}
\renewcommand{\braket}[2]{\left<#1| #2\right>}
\newcommand{\mc}[1]{\mathcal{#1}}
\newcommand{\ms}[1]{\mathscr{#1}}
\renewcommand{\md}[1]{\mathds{#1}}
\newcommand{\mb}[1]{\mathbb{#1}}
\newcommand{\tx}[1]{\mbox{#1}}
\newcommand{\mtx}[1]{\scriptstyle #1 \textstyle}
\newcommand{\ttx}[1]{\scriptscriptstyle #1 \textstyle}
\newcommand{\eqn}{equation}
\newcommand{\soln}{solution}
\newcommand{\curl}[1]{\vec{\nabla} \times \vec{#1}}
\newcommand{\grad}[1]{\vec{\nabla}#1}
\newcommand{\lap}[1]{\nabla^2 #1}
\newcommand{\mean}[1]{\langle #1\rangle}
\newcommand{\eref}[1]{Eq.~\eqref{#1}}
\newcommand{\fref}[1]{Fig.~\ref{#1}}
\newcommand{\vect}[1]{\boldsymbol{#1}}

\renewcommand{\d}[3]{\frac{d^{#3} #1}{d #2^{#3}}}
\renewcommand{\t}[1]{\tilde{#1}}
\renewcommand{\vec}[1]{\boldsymbol{#1}}
\renewcommand{\implies}{\Rightarrow}
\renewcommand{\div}[1]{\vec{\nabla} \cdot{#1}}

\DeclareMathOperator{\sech}{sech}
\DeclareMathOperator{\cosech}{cosech}
\DeclareMathOperator{\arcsec}{arcsec}
\DeclareMathOperator{\arccot}{arcCot}
\DeclareMathOperator{\arccsc}{arcCsc}
\DeclareMathOperator{\arccosh}{arcCosh}
\DeclareMathOperator{\arcsinh}{arcsinh}
\DeclareMathOperator{\arctanh}{arctanh}
\DeclareMathOperator{\arcsech}{arcsech}
\DeclareMathOperator{\arccsch}{arcCsch}
\DeclareMathOperator{\arccoth}{arcCoth}
\DeclareMathOperator{\Tr}{Tr}

% Margin/inline comment macros used across the original manuscripts.
% Rendered in color in draft mode; harmless to leave in for the thesis.
\makeatletter
\DeclareRobustCommand{\rc}[1]{\textcolor{olive}{#1}}
\DeclareRobustCommand{\rg}[1]{\textcolor{magenta}{#1}}
\DeclareRobustCommand{\rgc}[1]{\textcolor{magenta}{{[\bf #1]}}}
\DeclareRobustCommand{\rcomm}[1]{\textcolor{OliveGreen}{[#1]}}
\DeclareRobustCommand{\revise}[1]{\textcolor{blue}{#1}}
\makeatother
\newcommand{\hjc}[1]{\textcolor{blue}{#1}}

% Chapter-specific macros that appear inside figure or table captions. They
% have to live here rather than in their chapter, because \listoffigures and
% \listoftables typeset the .lof and .lot files in the front matter, long
% before the chapters that define them are read.
% Internal Tide-Eddy run names:
\newcommand{\op}[1]{\operatorname{#1}}
\newcommand{\roeff}{Ro_{\op{eff}}}
\newcommand{\lr}{\texttt{Ro-L}\xspace}
\newcommand{\ir}{\texttt{Ro-I}\xspace}
\newcommand{\hr}{\texttt{Ro-H}\xspace}
\newcommand{\lre}{\texttt{Ro-L8}\xspace}
\newcommand{\lro}{\texttt{Ro-L1}\xspace}
\newcommand{\ire}{\texttt{Ro-I8}\xspace}
\newcommand{\iro}{\texttt{Ro-I1}\xspace}
\newcommand{\hre}{\texttt{Ro-H8}\xspace}
\newcommand{\hro}{\texttt{Ro-H1}\xspace}
\newcommand{\filler}[1]{{\centering\bfseries #1\par}}
% Inertial Caustics particle labels:
\newcommand{\cs}{C~}
\newcommand{\su}{S~}
